\documentclass[11pt]{article}
\usepackage[margin=1in]{geometry}
\usepackage{amsmath, amssymb}
\usepackage{listings}
\usepackage{xcolor}

\lstset{
    basicstyle=\ttfamily\small,
    keywordstyle=\color{blue},
    commentstyle=\color{gray},
    stringstyle=\color{green!50!black},
    breaklines=true,
    frame=single
}
\usepackage{booktabs}
\usepackage{tabularx}
\usepackage{array}
\renewcommand{\arraystretch}{1.25}

\title{C++ para Programação Competitiva: STL, Heaps e Iteradores}
\date{}

\begin{document}

\maketitle

\section{Priority Queue e Heaps}

\subsection{Conceito de Heap}

Um heap é uma estrutura que mantém o elemento extremo acessível:
\begin{itemize}
    \item Max-heap: maior elemento no topo
    \item Min-heap: menor elemento no topo
\end{itemize}

Complexidade:
\[
\text{push} = O(\log n), \quad \text{pop} = O(\log n), \quad \text{top} = O(1)
\]

\subsection{Priority Queue padrão (max-heap)}

\begin{lstlisting}[language=C++]
priority_queue<int> pq;
pq.push(3);
pq.push(10);
pq.push(5);

cout << pq.top(); // 10
pq.pop();
\end{lstlisting}

\subsection{Min-heap}

\begin{lstlisting}[language=C++]
priority_queue<int, vector<int>, greater<int>> pq;
\end{lstlisting}

\subsection{Comparator customizado}

\begin{lstlisting}[language=C++]
auto cmp = [](pair<int,int> a, pair<int,int> b){
    return a.second > b.second;
};

priority_queue<pair<int,int>, vector<pair<int,int>>, decltype(cmp)> pq(cmp);
\end{lstlisting}

\subsection{Casos clássicos}

\begin{itemize}
    \item Dijkstra (min-heap)
    \item Top-K elementos
    \item Merge de listas ordenadas
    \item Scheduling (greedy)
\end{itemize}

\subsection{Limitações}

\begin{itemize}
    \item Não permite remover elementos arbitrários
    \item Não permite acesso interno
\end{itemize}

\section{Heaps manuais}

\begin{lstlisting}[language=C++]
vector<int> v = {3,1,5};

make_heap(v.begin(), v.end());
pop_heap(v.begin(), v.end());
v.pop_back();

v.push_back(10);
push_heap(v.begin(), v.end());
\end{lstlisting}

Min-heap:
\begin{lstlisting}[language=C++]
make_heap(v.begin(), v.end(), greater<int>());
\end{lstlisting}

\section{Iteradores}

\subsection{Definição}

Iteradores são ponteiros generalizados que percorrem containers.

\begin{lstlisting}[language=C++]
vector<int> v = {1,2,3};

for(auto it = v.begin(); it != v.end(); it++){
    cout << *it;
}
\end{lstlisting}

\subsection{Tipos de iteradores}

\begin{itemize}
    \item Random access: vector, deque
    \item Bidirectional: set, map
    \item Forward: unordered\_map
\end{itemize}

\subsection{Operações importantes}

\begin{lstlisting}[language=C++]
auto it = v.begin();
advance(it, k);
distance(a, b);
\end{lstlisting}

\subsection{Iteradores reversos}

\begin{lstlisting}[language=C++]
for(auto it = v.rbegin(); it != v.rend(); it++){
}
\end{lstlisting}

\subsection{Segurança}

Forma incorreta:
\begin{lstlisting}[language=C++]
v.erase(it); // invalida iterador
\end{lstlisting}

Forma correta:
\begin{lstlisting}[language=C++]
it = v.erase(it);
\end{lstlisting}

\section{Iteradores em set/map}

\begin{lstlisting}[language=C++]
set<int> s = {1,3,5};

auto it = s.lower_bound(3);
auto it2 = s.upper_bound(3);
\end{lstlisting}

\section{Comparação: priority\_queue vs set}

\begin{center}
\begin{tabular}{|c|c|c|c|}
\hline
Estrutura & Inserção & Remoção arbitrária & Ordenação \\
\hline
priority\_queue & $O(\log n)$ & Não & Parcial \\
set & $O(\log n)$ & Sim & Total \\
\hline
\end{tabular}
\end{center}

\section{Padrões com iteradores}

\subsection{Range-based loop}

\begin{lstlisting}[language=C++]
for(auto x : v)
\end{lstlisting}

\subsection{Loop com índice}

\begin{lstlisting}[language=C++]
for(int i = 0; i < v.size(); i++)
\end{lstlisting}

\section{Dicas avançadas}

\subsection{Lazy deletion em heap}

\begin{lstlisting}[language=C++]
unordered_map<int,int> removed;
\end{lstlisting}

\subsection{Heap com struct}

\begin{lstlisting}[language=C++]
struct Node{
    int dist, v;
    bool operator<(const Node& other) const{
        return dist > other.dist;
    }
};

priority_queue<Node> pq;
\end{lstlisting}

\section{Complexidade}

\begin{center}
\begin{tabular}{|c|c|c|c|}
\hline
Estrutura & Inserção & Remoção & Busca \\
\hline
vector & $O(1)$* & $O(n)$ & $O(n)$ \\
set/map & $O(\log n)$ & $O(\log n)$ & $O(\log n)$ \\
priority\_queue & $O(\log n)$ & $O(\log n)$ & topo $O(1)$ \\
\hline
\end{tabular}
\end{center}

\section{Exemplo 1: Dijkstra (priority\_queue)}

Problema: encontrar menor caminho em grafo ponderado.

\begin{lstlisting}[language=C++]
const int INF = 1e9;
vector<vector<pair<int,int>>> adj; // {vizinho, peso}

vector<int> dijkstra(int n, int src){
    vector<int> dist(n, INF);
    priority_queue<pair<int,int>, vector<pair<int,int>>, greater<>> pq;

    dist[src] = 0;
    pq.push({0, src});

    while(!pq.empty()){
        auto [d, u] = pq.top();
        pq.pop();

        if(d > dist[u]) continue;

        for(auto [v, w] : adj[u]){
            if(dist[v] > d + w){
                dist[v] = d + w;
                pq.push({dist[v], v});
            }
        }
    }
    return dist;
}
\end{lstlisting}

\textbf{Ideia:} sempre expandir o nó com menor distância atual.

\section{Exemplo 2: Top-K elementos}

Problema: manter os K maiores elementos.

\begin{lstlisting}[language=C++]
priority_queue<int, vector<int>, greater<int>> pq;

for(int x : v){
    pq.push(x);
    if(pq.size() > k) pq.pop();
}
\end{lstlisting}

\textbf{Ideia:} min-heap guarda apenas os K maiores.

\section{Exemplo 3: Sweep Line com multiset}

Problema: máximo número de intervalos ativos.

\begin{lstlisting}[language=C++]
vector<pair<int,int>> events;

for(auto [l, r] : intervals){
    events.push_back({l, 1});
    events.push_back({r, -1});
}

sort(events.begin(), events.end());

int cur = 0, ans = 0;
for(auto [x, type] : events){
    cur += type;
    ans = max(ans, cur);
}
\end{lstlisting}

\textbf{Ideia:} transformar intervalos em eventos.

\section{Exemplo 4: Sliding Window (two pointers)}

Problema: maior subarray com soma $\leq S$.

\begin{lstlisting}[language=C++]
int l = 0, sum = 0, ans = 0;

for(int r = 0; r < n; r++){
    sum += v[r];
    while(sum > S){
        sum -= v[l++];
    }
    ans = max(ans, r - l + 1);
}
\end{lstlisting}

\section{Exemplo 5: Uso de set (ordenado)}

Problema: encontrar menor elemento $\geq x$.

\begin{lstlisting}[language=C++]
set<int> s = {1,3,5,7};

auto it = s.lower_bound(4);
if(it != s.end()){
    cout << *it; // 5
}
\end{lstlisting}

\textbf{Ideia:} busca binária em estrutura dinâmica.

\section{Exemplo 6: Merge de K listas}

\begin{lstlisting}[language=C++]
priority_queue<
    pair<int,int>,
    vector<pair<int,int>>,
    greater<>
> pq;


\end{lstlisting}

\textbf{Ideia:} sempre pegar o menor entre os heads.

\section{Exemplo 7: Frequência com map}

\begin{lstlisting}[language=C++]
map<int,int> freq;

for(int x : v){
    freq[x]++;
}
\end{lstlisting}

\section{Exemplo 8: Binary Search}

\begin{lstlisting}[language=C++]
int l = 0, r = n-1;

while(l <= r){
    int m = (l + r) / 2;
    if(v[m] == x) return m;
    else if(v[m] < x) l = m + 1;
    else r = m - 1;
}
\end{lstlisting}

\section{Exemplo 9: Prefix Sum}

\begin{lstlisting}[language=C++]
vector<int> pref(n+1, 0);

for(int i = 0; i < n; i++){
    pref[i+1] = pref[i] + v[i];
}

// soma intervalo [l, r]
int sum = pref[r+1] - pref[l];
\end{lstlisting}

\section{Exemplo 10: DFS}

\begin{lstlisting}[language=C++]
vector<vector<int>> adj;
vector<bool> vis;

void dfs(int u){
    vis[u] = true;
    for(int v : adj[u]){
        if(!vis[v]) dfs(v);
    }
}
\end{lstlisting}

\section{Conclusão}

Esses padrões cobrem grande parte dos problemas de programação competitiva:
\begin{itemize}
    \item Grafos: Dijkstra, DFS/BFS
    \item Estruturas: heap, set, map
    \item Técnicas: two pointers, sweep line, prefix sum
    \item Busca: binary search
\end{itemize}

\section{Busca Binária e Funções Associadas (STL)}

\subsection{Pré-requisito}

Todas essas funções assumem que o container está ordenado:
\begin{lstlisting}[language=C++]
sort(v.begin(), v.end());
\end{lstlisting}

\subsection{binary\_search}

Verifica se um elemento existe.

\begin{lstlisting}[language=C++]
bool exists = binary_search(v.begin(), v.end(), x);
\end{lstlisting}

\[
\text{Retorna true se } x \in v, \text{ caso contrário false}
\]

Complexidade: \(O(\log n)\)

\bigskip

\subsection{lower\_bound}

Retorna um iterador para o \textbf{primeiro elemento $\geq x$}.

\begin{lstlisting}[language=C++]
auto it = lower_bound(v.begin(), v.end(), x);
\end{lstlisting}

Interpretação:
\[
\text{it} = \min \{ i \mid v[i] \geq x \}
\]

Exemplo:
\begin{lstlisting}[language=C++]
v = {1,2,4,4,5};

lower_bound(v.begin(), v.end(), 4) -> aponta para o primeiro 4
lower_bound(v.begin(), v.end(), 3) -> aponta para 4
\end{lstlisting}

\bigskip

\subsection{upper\_bound}

Retorna um iterador para o \textbf{primeiro elemento $> x$}.

\begin{lstlisting}[language=C++]
auto it = upper_bound(v.begin(), v.end(), x);
\end{lstlisting}

Interpretação:
\[
\text{it} = \min \{ i \mid v[i] > x \}
\]

Exemplo:
\begin{lstlisting}[language=C++]
v = {1,2,4,4,5};

upper_bound(v.begin(), v.end(), 4) -> aponta para 5
\end{lstlisting}

\bigskip

\subsection{Contagem de ocorrências}

Número de vezes que \(x\) aparece:

\begin{lstlisting}[language=C++]
int count = upper_bound(v.begin(), v.end(), x)
          - lower_bound(v.begin(), v.end(), x);
\end{lstlisting}

\[
\text{count} = |\{i : v[i] = x\}|
\]

\bigskip

\subsection{Busca binária manual}

\begin{lstlisting}[language=C++]
int l = 0, r = n-1;

while(l <= r){
    int m = l + (r - l) / 2;
    if(v[m] == x) return m;
    else if(v[m] < x) l = m + 1;
    else r = m - 1;
}
\end{lstlisting}

\bigskip

\subsection{Binary search em resposta (muito importante)}

Usado quando a resposta é monotônica.

\begin{lstlisting}[language=C++]
int l = 0, r = 1e9;

while(l < r){
    int m = (l + r) / 2;
    if(check(m)) r = m;
    else l = m + 1;
}
\end{lstlisting}

\[
\text{Objetivo: encontrar o menor } x \text{ tal que } check(x) = true
\]

\bigskip

\subsection{Uso em set/map}

\begin{lstlisting}[language=C++]
set<int> s = {1,3,5,7};

auto it = s.lower_bound(4); // aponta para 5
auto it2 = s.upper_bound(5); // aponta para 7
\end{lstlisting}

Complexidade: \(O(\log n)\)

\bigskip

\subsection{Resumo conceitual}

\begin{center}
\begin{tabular}{|c|c|}
\hline
Função & Significado \\
\hline
binary\_search & verifica existência \\
lower\_bound & primeiro $\geq x$ \\
upper\_bound & primeiro $> x$ \\
\hline
\end{tabular}
\end{center}

\bigskip

\subsection{Dicas práticas}

\begin{itemize}
    \item Sempre garantir que o vetor está ordenado
    \item Usar lower\_bound para encontrar posição de inserção
    \item Diferença de iteradores dá índice:
    \begin{lstlisting}[language=C++]
int idx = it - v.begin();
    \end{lstlisting}
    \item Evitar overflow:
    \[
    m = l + \frac{r - l}{2}
    \]
\end{itemize}

\section{Técnicas Avançadas de Busca e Estruturas}

\subsection{Ternary Search}

Usado para funções unimodais (cresce e depois decresce, ou vice-versa).

\textbf{Caso contínuo:}

\begin{lstlisting}[language=C++]
double ternary(double l, double r){
    for(int i = 0; i < 100; i++){
        double m1 = l + (r - l) / 3;
        double m2 = r - (r - l) / 3;

        if(f(m1) < f(m2)) l = m1;
        else r = m2;
    }
    return l;
}
\end{lstlisting}

\textbf{Caso discreto:}

\begin{lstlisting}[language=C++]
int ternary(int l, int r){
    while(r - l > 3){
        int m1 = l + (r - l) / 3;
        int m2 = r - (r - l) / 3;

        if(f(m1) < f(m2)) l = m1;
        else r = m2;
    }

    int ans = l;
    for(int i = l; i <= r; i++){
        if(f(i) < f(ans)) ans = i;
    }
    return ans;
}
\end{lstlisting}

\subsection{Busca binária em double}

Usada quando a resposta é contínua.

\begin{lstlisting}[language=C++]
double l = 0, r = 1e9;

for(int i = 0; i < 100; i++){
    double m = (l + r) / 2;
    if(check(m)) r = m;
    else l = m;
}
\end{lstlisting}

Critério alternativo:
\[
r - l < \varepsilon
\]

\subsection{Monotonic Stack}

Mantém elementos em ordem monotônica.

\textbf{Exemplo: próximo elemento maior à direita}

\begin{lstlisting}[language=C++]
vector<int> next_greater(n, -1);
stack<int> st;

for(int i = 0; i < n; i++){
    while(!st.empty() && v[st.top()] < v[i]){
        next_greater[st.top()] = v[i];
        st.pop();
    }
    st.push(i);
}
\end{lstlisting}

Complexidade: \(O(n)\)

\subsection{Monotonic Deque (Sliding Window)}

Usado para máximo/mínimo em janela.

\begin{lstlisting}[language=C++]
deque<int> dq;

for(int i = 0; i < n; i++){
    while(!dq.empty() && dq.front() <= i - k)
        dq.pop_front();

    while(!dq.empty() && v[dq.back()] <= v[i])
        dq.pop_back();

    dq.push_back(i);

    if(i >= k-1)
        cout << v[dq.front()] << " ";
}
\end{lstlisting}

\subsection{Resumo conceitual}

\begin{center}
\begin{tabular}{|c|c|}
\hline
Técnica & Uso \\
\hline
Binary Search & função monotônica \\
Ternary Search & função unimodal \\
Double Search & precisão contínua \\
Monotonic Stack & nearest greater/smaller \\
Monotonic Deque & sliding window max/min \\
\hline
\end{tabular}
\end{center}

\subsection{Dicas práticas}

\begin{itemize}
    \item Binary search: condição monotônica clara
    \item Ternary search: evitar se função não for unimodal
    \item Double search: usar número fixo de iterações
    \item Monotonic stack: cada elemento entra e sai no máximo uma vez
    \item Deque: mantém candidatos relevantes apenas
\end{itemize}

\section{Árvores, DFS e BFS}

\subsection{Definição de Árvore}

Uma árvore é um grafo não direcionado, conexo e sem ciclos.

Propriedades:
\begin{itemize}
    \item Possui $n$ vértices e $n-1$ arestas
    \item Existe exatamente um caminho entre quaisquer dois nós
    \item Pode ser enraizada (rooted tree)
\end{itemize}

Representação comum:
\begin{lstlisting}[language=C++]
vector<vector<int>> adj(n);
\end{lstlisting}

\subsection{DFS (Depth-First Search)}

Explora o grafo em profundidade.

\subsubsection{Implementação recursiva}

\begin{lstlisting}[language=C++]
vector<vector<int>> adj;
vector<bool> vis;

void dfs(int u){
    vis[u] = true;
    for(int v : adj[u]){
        if(!vis[v]){
            dfs(v);
        }
    }
}
\end{lstlisting}

\subsubsection{Complexidade}

\[
O(n + m)
\]

\subsubsection{Usos clássicos}

\begin{itemize}
    \item Encontrar componentes conexas
    \item Detectar ciclos
    \item Ordenação topológica
    \item DP em árvores
\end{itemize}

\subsection{BFS (Breadth-First Search)}

Explora o grafo em camadas (níveis).

\subsubsection{Implementação}

\begin{lstlisting}[language=C++]
vector<vector<int>> adj;
vector<int> dist(n, -1);

void bfs(int src){
    queue<int> q;
    q.push(src);
    dist[src] = 0;

    while(!q.empty()){
        int u = q.front(); q.pop();
        for(int v : adj[u]){
            if(dist[v] == -1){
                dist[v] = dist[u] + 1;
                q.push(v);
            }
        }
    }
}
\end{lstlisting}

\subsubsection{Complexidade}

\[
O(n + m)
\]

\subsubsection{Usos clássicos}

\begin{itemize}
    \item Menor caminho em grafos não ponderados
    \item Verificação de bipartição
    \item Expansão por níveis
\end{itemize}

\subsection{Diferença entre DFS e BFS}

\begin{center}
\begin{tabular}{|c|c|c|}
\hline
Característica & DFS & BFS \\
\hline
Estrutura & Stack / Recursão & Queue \\
Ordem & Profundidade & Níveis \\
Menor caminho & Não garante & Garante (não ponderado) \\
\hline
\end{tabular}
\end{center}

\subsection{DFS em árvore (com pai)}

\begin{lstlisting}[language=C++]
void dfs(int u, int p){
    for(int v : adj[u]){
        if(v != p){
            dfs(v, u);
        }
    }
}
\end{lstlisting}

\subsection{Altura e profundidade}

\begin{itemize}
    \item Profundidade: distância da raiz até o nó
    \item Altura: maior distância até uma folha
\end{itemize}

\subsection{Exemplo: calcular profundidade}

\begin{lstlisting}[language=C++]
vector<int> depth;

void dfs(int u, int p){
    for(int v : adj[u]){
        if(v != p){
            depth[v] = depth[u] + 1;
            dfs(v, u);
        }
    }
}
\end{lstlisting}

\subsection{Observações práticas}

\begin{itemize}
    \item Sempre marcar nós visitados para evitar loops
    \item Em árvores, usar parâmetro \texttt{parent} ao invés de \texttt{visited}
    \item BFS usa mais memória que DFS
\end{itemize}

\subsection{Lista de Adjacência (adj)}

A lista de adjacência é a forma padrão de representar grafos.

\subsubsection{Definição}

\begin{lstlisting}[language=C++]
vector<vector<int>> adj(n);
\end{lstlisting}

Interpretação:
\[
adj[u] = \{ \text{vizinhos de } u \}
\]

\subsubsection{Construção}

\textbf{Grafo não direcionado:}
\begin{lstlisting}[language=C++]
adj[u].push_back(v);
adj[v].push_back(u);
\end{lstlisting}

\textbf{Grafo direcionado:}
\begin{lstlisting}[language=C++]
adj[u].push_back(v);
\end{lstlisting}

\subsubsection{Grafo com pesos}

\begin{lstlisting}[language=C++]
vector<vector<pair<int,int>>> adj(n);
\end{lstlisting}

\[
adj[u] = \{ (v, w) \}
\]

\subsubsection{Iteração}

\begin{lstlisting}[language=C++]
for(int v : adj[u]) {
    // processa vizinho
}
\end{lstlisting}

Com pesos:
\begin{lstlisting}[language=C++]
for(auto [v, w] : adj[u]) {
    // usa v e w
}
\end{lstlisting}

\subsubsection{Complexidade}

\begin{itemize}
    \item Memória: $O(n + m)$
    \item Iteração DFS/BFS: $O(n + m)$
\end{itemize}

\subsubsection{Comparação com matriz}

\begin{center}
\begin{tabular}{|c|c|c|}
\hline
Estrutura & Memória & Iteração \\
\hline
Matriz & $O(n^2)$ & Lenta \\
Lista (adj) & $O(n+m)$ & Rápida \\
\hline
\end{tabular}
\end{center}

\section{Problemas Clássicos com Lista de Adjacência (adj)}

\subsection{1. Componentes Conexas (DFS)}

\textbf{Problema:} contar quantos componentes existem em um grafo.

\textbf{Ideia:} rodar DFS a partir de cada nó não visitado.

\begin{lstlisting}[language=C++]
vector<vector<int>> adj;
vector<bool> vis;

void dfs(int u){
    vis[u] = true;
    for(int v : adj[u]){
        if(!vis[v]) dfs(v);
    }
}

int count_components(int n){
    vis.assign(n, false);
    int count = 0;

    for(int i = 0; i < n; i++){
        if(!vis[i]){
            dfs(i);
            count++;
        }
    }
    return count;
}
\end{lstlisting}

---

\subsection{2. Menor Caminho (BFS)}

\textbf{Problema:} menor distância em grafo não ponderado.

\begin{lstlisting}[language=C++]
vector<int> dist(n, -1);

void bfs(int src){
    queue<int> q;
    q.push(src);
    dist[src] = 0;

    while(!q.empty()){
        int u = q.front(); q.pop();
        for(int v : adj[u]){
            if(dist[v] == -1){
                dist[v] = dist[u] + 1;
                q.push(v);
            }
        }
    }
}
\end{lstlisting}

---

\subsection{3. Detectar Ciclo (DFS)}

\textbf{Problema:} verificar se um grafo não direcionado tem ciclo.

\begin{lstlisting}[language=C++]
bool has_cycle(int u, int parent){
    vis[u] = true;
    for(int v : adj[u]){
        if(!vis[v]){
            if(has_cycle(v, u)) return true;
        } else if(v != parent){
            return true;
        }
    }
    return false;
}
\end{lstlisting}

---

\subsection{4. Árvore: calcular profundidade}

\textbf{Problema:} depth de cada nó.

\begin{lstlisting}[language=C++]
vector<int> depth;

void dfs(int u, int p){
    for(int v : adj[u]){
        if(v != p){
            depth[v] = depth[u] + 1;
            dfs(v, u);
        }
    }
}
\end{lstlisting}

---

\subsection{5. Diâmetro de Árvore}

\textbf{Problema:} maior distância entre dois nós.

\textbf{Ideia:} BFS duas vezes.

\begin{lstlisting}[language=C++]
int bfs_far(int src){
    vector<int> dist(n, -1);
    queue<int> q;

    q.push(src);
    dist[src] = 0;

    int last = src;

    while(!q.empty()){
        int u = q.front(); q.pop();
        last = u;
        for(int v : adj[u]){
            if(dist[v] == -1){
                dist[v] = dist[u] + 1;
                q.push(v);
            }
        }
    }
    return last;
}
\end{lstlisting}

---

\subsection{6. Bipartição (BFS)}

\textbf{Problema:} verificar se o grafo é bipartido.

\begin{lstlisting}[language=C++]
vector<int> color(n, -1);

bool is_bipartite(int src){
    queue<int> q;
    q.push(src);
    color[src] = 0;

    while(!q.empty()){
        int u = q.front(); q.pop();
        for(int v : adj[u]){
            if(color[v] == -1){
                color[v] = color[u] ^ 1;
                q.push(v);
            } else if(color[v] == color[u]){
                return false;
            }
        }
    }
    return true;
}
\end{lstlisting}

---

\subsection{7. Ordenação Topológica (DAG)}

\textbf{Problema:} ordenar grafo direcionado acíclico.

\begin{lstlisting}[language=C++]
vector<int> indeg(n, 0);

for(int u = 0; u < n; u++){
    for(int v : adj[u]){
        indeg[v]++;
    }
}

queue<int> q;
for(int i = 0; i < n; i++){
    if(indeg[i] == 0) q.push(i);
}

vector<int> topo;

while(!q.empty()){
    int u = q.front(); q.pop();
    topo.push_back(u);

    for(int v : adj[u]){
        if(--indeg[v] == 0){
            q.push(v);
        }
    }
}
\end{lstlisting}

---

\subsection{Resumo de padrões}

\begin{center}
\begin{tabular}{|c|c|}
\hline
Problema & Técnica \\
\hline
Componentes & DFS \\
Menor caminho & BFS \\
Ciclo & DFS \\
Árvore & DFS com pai \\
Diâmetro & BFS x2 \\
Bipartição & BFS + cores \\
Topológico & BFS (Kahn) \\
\hline
\end{tabular}
\end{center}

\begin{table}[h!]
\centering
\small
\begin{tabularx}{\textwidth}{>{\bfseries}p{2.4cm} X X p{2.5cm} p{2.2cm}}
\toprule
Algorithm & Main idea & Best use case & Data structure & Complexity \\
\midrule
Boruvka &
Each connected component chooses its minimum outgoing edge; all safe edges are added simultaneously. &
Minimum spanning tree when parallelization is useful or when we want to merge many components at once. &
Connected components, labels, Union-Find possible. &
$O(E \log V)$ \\
\midrule
Prim-Jarnik &
Grow one tree from an initial vertex; repeatedly add the cheapest edge leaving the current tree. &
Minimum spanning tree on connected weighted graphs, especially dense graphs or when using priority queues. &
Priority queue; Fibonacci heap for optimal bound. &
$O(E + V \log V)$ with Fibonacci heap \\
\midrule
Kruskal &
Sort edges by increasing weight; add an edge if it connects two different components. &
Minimum spanning tree, especially for sparse graphs. Very simple and practical. &
Union-Find / Disjoint Set Union. &
$O(E \log V)$ \\
\bottomrule
\end{tabularx}
\caption{Summary of minimum spanning tree algorithms.}
\end{table}

\begin{table}[h!]
\centering
\small
\begin{tabularx}{\textwidth}{>{\bfseries}p{2.8cm} X X X}
\toprule
Algorithm & Graph type & Edge weights allowed & Main characteristic \\
\midrule
BFS &
Unweighted graph, or graph where every edge has the same weight. &
All weights equal, usually $1$. &
Finds shortest paths by exploring vertices layer by layer from the source. \\
\midrule
DAG shortest paths &
Directed acyclic graph. &
Can allow negative weights. &
Processes vertices in topological order; works because there are no directed cycles. \\
\midrule
Dijkstra &
Directed or undirected weighted graph. &
Requires nonnegative weights. &
Greedy algorithm: repeatedly finalizes the vertex with smallest tentative distance. \\
\midrule
Bellman-Ford &
Directed or undirected weighted graph. &
Allows negative weights. &
Repeatedly relaxes all edges; can detect negative cycles. \\
\midrule
Johnson &
Directed weighted graph. &
Allows negative weights, but no negative cycles. &
Solves all-pairs shortest paths by reweighting edges, then running Dijkstra from every vertex. \\
\bottomrule
\end{tabularx}
\caption{Summary of shortest path algorithms.}
\end{table}

\section{Essential Graph Algorithms in C++}

\subsection{Graph Representation}

We use an adjacency list. For weighted graphs, each edge is stored as a pair
\texttt{(neighbor, weight)}.

\begin{lstlisting}
#include <bits/stdc++.h>
using namespace std;

const long long INF = 4e18;

struct Edge {
    int u, v;
    long long w;
};

using Graph = vector<vector<pair<int, long long>>>;
// graph[u] contains pairs (v, w), meaning edge u -> v with weight w.
\end{lstlisting}

\subsection{Breadth-First Search for Unweighted Shortest Paths}

BFS computes shortest paths when all edges have the same weight, usually weight 1.

\begin{lstlisting}
vector<int> bfs_shortest_paths(const vector<vector<int>>& graph, int source) {
    int n = graph.size();

    vector<int> dist(n, -1);
    vector<int> parent(n, -1);
    queue<int> q;

    dist[source] = 0;
    q.push(source);

    while (!q.empty()) {
        int u = q.front();
        q.pop();

        for (int v : graph[u]) {
            if (dist[v] == -1) {
                dist[v] = dist[u] + 1;
                parent[v] = u;
                q.push(v);
            }
        }
    }

    return dist;
}
\end{lstlisting}

Complexity: \texttt{O(V + E)}.

\subsection{Shortest Paths in a DAG}

For a directed acyclic graph, we process vertices in topological order and relax
all outgoing edges.

\begin{lstlisting}
void dfs_topological(
    int u,
    const Graph& graph,
    vector<bool>& visited,
    vector<int>& order
) {
    visited[u] = true;

    for (auto [v, w] : graph[u]) {
        if (!visited[v]) {
            dfs_topological(v, graph, visited, order);
        }
    }

    order.push_back(u);
}

vector<long long> dag_shortest_paths(const Graph& graph, int source) {
    int n = graph.size();

    vector<bool> visited(n, false);
    vector<int> order;

    for (int u = 0; u < n; u++) {
        if (!visited[u]) {
            dfs_topological(u, graph, visited, order);
        }
    }

    reverse(order.begin(), order.end());

    vector<long long> dist(n, INF);
    dist[source] = 0;

    for (int u : order) {
        if (dist[u] == INF) continue;

        for (auto [v, w] : graph[u]) {
            if (dist[u] + w < dist[v]) {
                dist[v] = dist[u] + w;
            }
        }
    }

    return dist;
}
\end{lstlisting}

Complexity: \texttt{O(V + E)}.

\subsection{Dijkstra's Algorithm}

Dijkstra computes shortest paths from one source when all edge weights are
nonnegative.

\begin{lstlisting}
vector<long long> dijkstra(const Graph& graph, int source) {
    int n = graph.size();

    vector<long long> dist(n, INF);
    dist[source] = 0;

    priority_queue<
        pair<long long, int>,
        vector<pair<long long, int>>,
        greater<pair<long long, int>>
    > pq;

    pq.push({0, source});

    while (!pq.empty()) {
        auto [du, u] = pq.top();
        pq.pop();

        if (du != dist[u]) continue;

        for (auto [v, w] : graph[u]) {
            if (dist[u] + w < dist[v]) {
                dist[v] = dist[u] + w;
                pq.push({dist[v], v});
            }
        }
    }

    return dist;
}
\end{lstlisting}

Complexity with binary heap: \texttt{O(E log V)}.

\subsection{Bellman-Ford Algorithm}

Bellman-Ford allows negative weights and detects negative cycles reachable from
the source.

\begin{lstlisting}
pair<vector<long long>, bool> bellman_ford(
    int n,
    const vector<Edge>& edges,
    int source
) {
    vector<long long> dist(n, INF);
    dist[source] = 0;

    for (int i = 0; i < n - 1; i++) {
        bool changed = false;

        for (const Edge& e : edges) {
            if (dist[e.u] == INF) continue;

            if (dist[e.u] + e.w < dist[e.v]) {
                dist[e.v] = dist[e.u] + e.w;
                changed = true;
            }
        }

        if (!changed) break;
    }

    bool has_negative_cycle = false;

    for (const Edge& e : edges) {
        if (dist[e.u] == INF) continue;

        if (dist[e.u] + e.w < dist[e.v]) {
            has_negative_cycle = true;
            break;
        }
    }

    return {dist, has_negative_cycle};
}
\end{lstlisting}

Complexity: \texttt{O(VE)}.

\subsection{Union-Find / Disjoint Set Union}

Union-Find is used by Kruskal's algorithm to test whether adding an edge creates
a cycle.

\begin{lstlisting}
struct DSU {
    vector<int> parent;
    vector<int> size;

    DSU(int n) {
        parent.resize(n);
        size.assign(n, 1);

        for (int i = 0; i < n; i++) {
            parent[i] = i;
        }
    }

    int find(int x) {
        if (parent[x] == x) return x;
        return parent[x] = find(parent[x]);
    }

    bool unite(int a, int b) {
        a = find(a);
        b = find(b);

        if (a == b) return false;

        if (size[a] < size[b]) {
            swap(a, b);
        }

        parent[b] = a;
        size[a] += size[b];

        return true;
    }
};
\end{lstlisting}

Amortized complexity per operation: almost constant.

\subsection{Kruskal's Algorithm}

Kruskal computes a minimum spanning tree by sorting edges and adding the
cheapest edges that do not create cycles.

\begin{lstlisting}
long long kruskal(int n, vector<Edge>& edges) {
    sort(edges.begin(), edges.end(), [](const Edge& a, const Edge& b) {
        return a.w < b.w;
    });

    DSU dsu(n);

    long long total_weight = 0;
    int edges_used = 0;

    for (const Edge& e : edges) {
        if (dsu.unite(e.u, e.v)) {
            total_weight += e.w;
            edges_used++;
        }
    }

    if (edges_used != n - 1) {
        throw runtime_error("Graph is not connected");
    }

    return total_weight;
}
\end{lstlisting}

Complexity: \texttt{O(E log V)}.

\subsection{Prim's Algorithm}

Prim computes a minimum spanning tree by growing one tree and always adding the
cheapest edge leaving the current tree.

\begin{lstlisting}
long long prim(const Graph& graph) {
    int n = graph.size();

    vector<bool> in_mst(n, false);
    vector<long long> min_edge(n, INF);

    priority_queue<
        pair<long long, int>,
        vector<pair<long long, int>>,
        greater<pair<long long, int>>
    > pq;

    min_edge[0] = 0;
    pq.push({0, 0});

    long long total_weight = 0;
    int vertices_used = 0;

    while (!pq.empty()) {
        auto [weight, u] = pq.top();
        pq.pop();

        if (in_mst[u]) continue;

        in_mst[u] = true;
        total_weight += weight;
        vertices_used++;

        for (auto [v, w] : graph[u]) {
            if (!in_mst[v] && w < min_edge[v]) {
                min_edge[v] = w;
                pq.push({w, v});
            }
        }
    }

    if (vertices_used != n) {
        throw runtime_error("Graph is not connected");
    }

    return total_weight;
}
\end{lstlisting}

Complexity with binary heap: \texttt{O(E log V)}.

\subsection{Johnson's Algorithm: Core Reweighting Idea}

Johnson's algorithm is used for all-pairs shortest paths with negative edges but
no negative cycles. The implementation below gives the essential structure.

\begin{lstlisting}
vector<vector<long long>> johnson(int n, vector<Edge> edges) {
    int artificial_source = n;

    vector<Edge> extended_edges = edges;

    for (int v = 0; v < n; v++) {
        extended_edges.push_back({artificial_source, v, 0});
    }

    auto [h, has_negative_cycle] =
        bellman_ford(n + 1, extended_edges, artificial_source);

    if (has_negative_cycle) {
        throw runtime_error("Negative cycle detected");
    }

    Graph reweighted_graph(n);

    for (const Edge& e : edges) {
        long long new_weight = e.w + h[e.u] - h[e.v];
        reweighted_graph[e.u].push_back({e.v, new_weight});
    }

    vector<vector<long long>> dist(n, vector<long long>(n, INF));

    for (int s = 0; s < n; s++) {
        vector<long long> d = dijkstra(reweighted_graph, s);

        for (int v = 0; v < n; v++) {
            if (d[v] < INF) {
                dist[s][v] = d[v] - h[s] + h[v];
            }
        }
    }

    return dist;
}
\end{lstlisting}

Complexity with binary heap: approximately \texttt{O(VE log V)}.
With Fibonacci heaps, the theoretical bound improves.

\subsection{Iterative DFS Using a Stack}

DFS can be implemented without recursion by using an explicit stack. This is
useful when the graph is large and recursive DFS may cause stack overflow.

\begin{lstlisting}
vector<int> dfs_iterative(const vector<vector<int>>& graph, int source) {
    int n = graph.size();

    vector<bool> visited(n, false);
    vector<int> order;

    stack<int> st;
    st.push(source);

    while (!st.empty()) {
        int u = st.top();
        st.pop();

        if (visited[u]) continue;

        visited[u] = true;
        order.push_back(u);

        for (int v : graph[u]) {
            if (!visited[v]) {
                st.push(v);
            }
        }
    }

    return order;
}
\end{lstlisting}

Complexity: \texttt{O(V + E)}.

\section{Important String Algorithms in C++}

\subsection{Naive String Matching}

The naive algorithm tries every possible starting position in the text and
compares the pattern character by character.

\begin{lstlisting}
#include <bits/stdc++.h>
using namespace std;

vector<int> naive_match(const string& text, const string& pattern) {
    int n = text.size();
    int m = pattern.size();

    vector<int> occurrences;

    for (int i = 0; i + m <= n; i++) {
        bool ok = true;

        for (int j = 0; j < m; j++) {
            if (text[i + j] != pattern[j]) {
                ok = false;
                break;
            }
        }

        if (ok) {
            occurrences.push_back(i);
        }
    }

    return occurrences;
}
\end{lstlisting}

Complexity: $O(nm)$ in the worst case.

\subsection{Polynomial Rolling Hash}

Polynomial hashing allows us to compare substrings in constant time after
linear preprocessing.

\begin{lstlisting}
struct RollingHash {
    static const long long MOD = 1000000007;
    static const long long BASE = 911382323;

    vector<long long> h;
    vector<long long> power;

    RollingHash(const string& s) {
        int n = s.size();

        h.assign(n + 1, 0);
        power.assign(n + 1, 1);

        for (int i = 0; i < n; i++) {
            h[i + 1] = (h[i] * BASE + s[i]) % MOD;
            power[i + 1] = (power[i] * BASE) % MOD;
        }
    }

    long long get_hash(int l, int r) {
        // Returns hash of s[l ... r - 1]
        long long res = h[r] - (h[l] * power[r - l]) % MOD;

        if (res < 0) {
            res += MOD;
        }

        return res;
    }
};
\end{lstlisting}

Preprocessing: $O(n)$.

Substring hash query: $O(1)$.

\subsection{Rabin-Karp String Matching}

Rabin-Karp uses rolling hash to compare the pattern with every substring of
the text of the same length.

\begin{lstlisting}
vector<int> rabin_karp(const string& text, const string& pattern) {
    int n = text.size();
    int m = pattern.size();

    vector<int> occurrences;

    if (m > n) {
        return occurrences;
    }

    RollingHash text_hash(text);
    RollingHash pattern_hash(pattern);

    long long target = pattern_hash.get_hash(0, m);

    for (int i = 0; i + m <= n; i++) {
        if (text_hash.get_hash(i, i + m) == target) {
            // Optional verification to avoid hash collision errors.
            if (text.substr(i, m) == pattern) {
                occurrences.push_back(i);
            }
        }
    }

    return occurrences;
}
\end{lstlisting}

Expected complexity: $O(n + m)$ if collisions are rare.

Worst case with many collisions: $O(nm)$ if we verify every collision.

\subsection{Prefix Function for KMP}

The prefix function $\pi[i]$ is the length of the longest proper prefix of
\texttt{pattern[0 ... i]} that is also a suffix of \texttt{pattern[0 ... i]}.

\begin{lstlisting}
vector<int> compute_pi(const string& pattern) {
    int m = pattern.size();

    vector<int> pi(m, 0);

    for (int i = 1; i < m; i++) {
        int j = pi[i - 1];

        while (j > 0 && pattern[i] != pattern[j]) {
            j = pi[j - 1];
        }

        if (pattern[i] == pattern[j]) {
            j++;
        }

        pi[i] = j;
    }

    return pi;
}
\end{lstlisting}

Complexity: $O(m)$.

\subsection{Knuth-Morris-Pratt Algorithm}

KMP finds all occurrences of a pattern in a text in linear time.

\begin{lstlisting}
vector<int> kmp_match(const string& text, const string& pattern) {
    int n = text.size();
    int m = pattern.size();

    vector<int> occurrences;

    if (m == 0) {
        return occurrences;
    }

    vector<int> pi = compute_pi(pattern);

    int j = 0;

    for (int i = 0; i < n; i++) {
        while (j > 0 && text[i] != pattern[j]) {
            j = pi[j - 1];
        }

        if (text[i] == pattern[j]) {
            j++;
        }

        if (j == m) {
            occurrences.push_back(i - m + 1);

            // Allows overlapping occurrences.
            j = pi[j - 1];
        }
    }

    return occurrences;
}
\end{lstlisting}

Complexity: $O(n + m)$.

\subsection{Trie}

A trie stores a set of strings by sharing common prefixes.

\begin{lstlisting}
struct TrieNode {
    TrieNode* child[26];
    bool is_end;

    TrieNode() {
        for (int i = 0; i < 26; i++) {
            child[i] = nullptr;
        }

        is_end = false;
    }
};

struct Trie {
    TrieNode* root;

    Trie() {
        root = new TrieNode();
    }

    void insert(const string& s) {
        TrieNode* cur = root;

        for (char c : s) {
            int x = c - 'a';

            if (cur->child[x] == nullptr) {
                cur->child[x] = new TrieNode();
            }

            cur = cur->child[x];
        }

        cur->is_end = true;
    }

    bool contains(const string& s) {
        TrieNode* cur = root;

        for (char c : s) {
            int x = c - 'a';

            if (cur->child[x] == nullptr) {
                return false;
            }

            cur = cur->child[x];
        }

        return cur->is_end;
    }

    bool starts_with(const string& prefix) {
        TrieNode* cur = root;

        for (char c : prefix) {
            int x = c - 'a';

            if (cur->child[x] == nullptr) {
                return false;
            }

            cur = cur->child[x];
        }

        return true;
    }
};
\end{lstlisting}

Insertion/search complexity: $O(L)$, where $L$ is the length of the word.

\subsection{Suffix Array}

A suffix array stores the starting positions of all suffixes of a string in
lexicographic order.

\begin{lstlisting}
vector<int> build_suffix_array(const string& s) {
    int n = s.size();

    vector<int> sa(n);
    vector<int> rank(n);
    vector<int> tmp(n);

    for (int i = 0; i < n; i++) {
        sa[i] = i;
        rank[i] = s[i];
    }

    for (int k = 1; k < n; k *= 2) {
        auto cmp = [&](int i, int j) {
            if (rank[i] != rank[j]) {
                return rank[i] < rank[j];
            }

            int ri = (i + k < n ? rank[i + k] : -1);
            int rj = (j + k < n ? rank[j + k] : -1);

            return ri < rj;
        };

        sort(sa.begin(), sa.end(), cmp);

        tmp[sa[0]] = 0;

        for (int i = 1; i < n; i++) {
            tmp[sa[i]] = tmp[sa[i - 1]] + cmp(sa[i - 1], sa[i]);
        }

        for (int i = 0; i < n; i++) {
            rank[i] = tmp[i];
        }

        if (rank[sa[n - 1]] == n - 1) {
            break;
        }
    }

    return sa;
}
\end{lstlisting}

Complexity: $O(n \log^2 n)$ with comparison sorting.

\subsection{Pattern Search with Suffix Array}

Since the suffixes are sorted, we can binary search for a pattern.

\begin{lstlisting}
bool starts_with_pattern(const string& s, int pos, const string& pattern) {
    int n = s.size();
    int m = pattern.size();

    if (pos + m > n) {
        return false;
    }

    for (int i = 0; i < m; i++) {
        if (s[pos + i] != pattern[i]) {
            return false;
        }
    }

    return true;
}

bool suffix_less_than_pattern(const string& s, int pos, const string& pattern) {
    int n = s.size();
    int m = pattern.size();

    int i = 0;

    while (pos + i < n && i < m) {
        if (s[pos + i] != pattern[i]) {
            return s[pos + i] < pattern[i];
        }

        i++;
    }

    return i == m ? false : true;
}

vector<int> suffix_array_match(
    const string& s,
    const vector<int>& sa,
    const string& pattern
) {
    int n = s.size();

    int left = 0;
    int right = n;

    while (left < right) {
        int mid = (left + right) / 2;

        if (suffix_less_than_pattern(s, sa[mid], pattern)) {
            left = mid + 1;
        } else {
            right = mid;
        }
    }

    vector<int> occurrences;

    for (int i = left; i < n; i++) {
        if (starts_with_pattern(s, sa[i], pattern)) {
            occurrences.push_back(sa[i]);
        } else {
            break;
        }
    }

    sort(occurrences.begin(), occurrences.end());

    return occurrences;
}
\end{lstlisting}

Basic query complexity: $O(m \log n + km)$, where $k$ is the number of
occurrences.

\subsection{LCP Array: Kasai's Algorithm}

The LCP array stores the longest common prefix between consecutive suffixes in
the suffix array.

\begin{lstlisting}
vector<int> build_lcp(const string& s, const vector<int>& sa) {
    int n = s.size();

    vector<int> rank(n);
    vector<int> lcp(n - 1);

    for (int i = 0; i < n; i++) {
        rank[sa[i]] = i;
    }

    int h = 0;

    for (int i = 0; i < n; i++) {
        int r = rank[i];

        if (r == n - 1) {
            h = 0;
            continue;
        }

        int j = sa[r + 1];

        while (i + h < n && j + h < n && s[i + h] == s[j + h]) {
            h++;
        }

        lcp[r] = h;

        if (h > 0) {
            h--;
        }
    }

    return lcp;
}
\end{lstlisting}

Complexity: $O(n)$ after the suffix array is built.

\subsection{Longest Palindromic Subsequence}

The longest palindromic subsequence of a string can be computed by dynamic
programming.

\begin{lstlisting}
int longest_palindromic_subsequence(const string& s) {
    int n = s.size();

    vector<vector<int>> dp(n, vector<int>(n, 0));

    for (int i = 0; i < n; i++) {
        dp[i][i] = 1;
    }

    for (int len = 2; len <= n; len++) {
        for (int l = 0; l + len <= n; l++) {
            int r = l + len - 1;

            if (s[l] == s[r]) {
                dp[l][r] = 2;

                if (l + 1 <= r - 1) {
                    dp[l][r] += dp[l + 1][r - 1];
                }
            } else {
                dp[l][r] = max(dp[l + 1][r], dp[l][r - 1]);
            }
        }
    }

    return dp[0][n - 1];
}
\end{lstlisting}

Complexity: $O(n^2)$.

\subsection{General Backtracking Framework}

Backtracking is used when we build a solution step by step and abandon a
partial solution as soon as it cannot lead to a valid complete solution.

The general structure is:

\begin{lstlisting}
void backtrack(State& state) {
    if (is_complete(state)) {
        process_solution(state);
        return;
    }

    for (Choice choice : possible_choices(state)) {
        if (is_valid(choice, state)) {
            apply_choice(choice, state);

            backtrack(state);

            undo_choice(choice, state);
        }
    }
}
\end{lstlisting}

\subsection{Backtracking Template in C++}

\begin{lstlisting}
#include <bits/stdc++.h>
using namespace std;

struct State {
    vector<int> solution;
};

bool is_complete(const State& state) {
    // Return true if the current partial solution is complete.
    return false;
}

vector<int> possible_choices(const State& state) {
    // Return the list of choices that we can try from this state.
    return {};
}

bool is_valid(int choice, const State& state) {
    // Return true if adding this choice keeps the state valid.
    return true;
}

void apply_choice(int choice, State& state) {
    // Add the choice to the current solution.
    state.solution.push_back(choice);
}

void undo_choice(int choice, State& state) {
    // Remove the choice and restore the previous state.
    state.solution.pop_back();
}

void process_solution(const State& state) {
    // Use or store the complete solution.
    for (int x : state.solution) {
        cout << x << " ";
    }
    cout << "\n";
}

void backtrack(State& state) {
    if (is_complete(state)) {
        process_solution(state);
        return;
    }

    for (int choice : possible_choices(state)) {
        if (is_valid(choice, state)) {
            apply_choice(choice, state);

            backtrack(state);

            undo_choice(choice, state);
        }
    }
}
\end{lstlisting}

\subsection{Example 1: Generate All Subsets}

Given an array, we decide for each element whether to take it or not.

\begin{lstlisting}
#include <bits/stdc++.h>
using namespace std;

vector<int> a;
vector<int> current;
vector<vector<int>> all_subsets;

void backtrack_subsets(int i) {
    if (i == (int)a.size()) {
        all_subsets.push_back(current);
        return;
    }

    // Choice 1: do not take a[i]
    backtrack_subsets(i + 1);

    // Choice 2: take a[i]
    current.push_back(a[i]);
    backtrack_subsets(i + 1);
    current.pop_back();
}

int main() {
    a = {1, 2, 3};

    backtrack_subsets(0);

    for (auto subset : all_subsets) {
        cout << "{ ";
        for (int x : subset) {
            cout << x << " ";
        }
        cout << "}\n";
    }

    return 0;
}
\end{lstlisting}

Complexity: $O(2^n)$ subsets.

\subsection{Example 2: Generate All Permutations}

We build the permutation position by position. At each step, we choose an
unused element.

\begin{lstlisting}
#include <bits/stdc++.h>
using namespace std;

int n;
vector<int> perm;
vector<bool> used;

void backtrack_permutations() {
    if ((int)perm.size() == n) {
        for (int x : perm) {
            cout << x << " ";
        }
        cout << "\n";
        return;
    }

    for (int x = 1; x <= n; x++) {
        if (!used[x]) {
            used[x] = true;
            perm.push_back(x);

            backtrack_permutations();

            perm.pop_back();
            used[x] = false;
        }
    }
}

int main() {
    n = 3;
    used.assign(n + 1, false);

    backtrack_permutations();

    return 0;
}
\end{lstlisting}

Complexity: $O(n!)$ permutations.

\subsection{Example 3: Generate All Combinations of Size $k$}

We choose $k$ elements among $\{1,\dots,n\}$.

\begin{lstlisting}
#include <bits/stdc++.h>
using namespace std;

int n, k;
vector<int> comb;

void backtrack_combinations(int start) {
    if ((int)comb.size() == k) {
        for (int x : comb) {
            cout << x << " ";
        }
        cout << "\n";
        return;
    }

    for (int x = start; x <= n; x++) {
        comb.push_back(x);

        backtrack_combinations(x + 1);

        comb.pop_back();
    }
}

int main() {
    n = 5;
    k = 3;

    backtrack_combinations(1);

    return 0;
}
\end{lstlisting}

Complexity: $O\binom{n}{k}$ generated solutions.

\subsection{Example 4: N-Queens Backtracking}

Place $n$ queens on an $n \times n$ chessboard so that no two queens attack
each other.

\begin{lstlisting}
#include <bits/stdc++.h>
using namespace std;

int n;
vector<int> queen_col;

// used_col[c] tells whether column c already has a queen.
// used_diag1[r + c] corresponds to one diagonal direction.
// used_diag2[r - c + n - 1] corresponds to the other diagonal direction.
vector<bool> used_col;
vector<bool> used_diag1;
vector<bool> used_diag2;

void print_solution() {
    for (int r = 0; r < n; r++) {
        for (int c = 0; c < n; c++) {
            if (queen_col[r] == c) {
                cout << "Q";
            } else {
                cout << ".";
            }
        }
        cout << "\n";
    }
    cout << "\n";
}

void backtrack_queens(int row) {
    if (row == n) {
        print_solution();
        return;
    }

    for (int col = 0; col < n; col++) {
        int d1 = row + col;
        int d2 = row - col + n - 1;

        if (used_col[col] || used_diag1[d1] || used_diag2[d2]) {
            continue;
        }

        queen_col[row] = col;
        used_col[col] = true;
        used_diag1[d1] = true;
        used_diag2[d2] = true;

        backtrack_queens(row + 1);

        queen_col[row] = -1;
        used_col[col] = false;
        used_diag1[d1] = false;
        used_diag2[d2] = false;
    }
}

int main() {
    n = 8;

    queen_col.assign(n, -1);
    used_col.assign(n, false);
    used_diag1.assign(2 * n - 1, false);
    used_diag2.assign(2 * n - 1, false);

    backtrack_queens(0);

    return 0;
}
\end{lstlisting}

\end{document}